% VI/05 exercise hints. Problems and Challenges carry guided hints in their dossiers.

\subsection*{Hint to Exercise VI.05.1}
Compute the three coordinate rings and ask whether each is a domain.

\subsection*{Hint to Exercise VI.05.2}
Every nonempty open subset of an irreducible space is dense.

\subsection*{Hint to Exercise VI.05.3}
The function \(f-g\) vanishes on a dense subset, and its zero set is closed.

\subsection*{Hint to Exercise VI.05.4}
The affine line is irreducible.  Use continuity of the parametrization and then its image.

\subsection*{Hint to Exercise VI.05.5}
Factor the defining polynomial into four distinct linear factors.  All four lines contain
the origin.

\subsection*{Hint to Exercise VI.05.6}
The equation \(xy=0\) says \(x=0\) or \(y=0\).

\subsection*{Hint to Exercise VI.05.7}
For idempotents, use connectedness of the union of the two axes, or restrict the idempotent
to each irreducible component and compare values at the origin.

\subsection*{Hint to Exercise VI.05.8}
The ideals \((x)\) and \((x-1)\) are comaximal.

\subsection*{Hint to Exercise VI.05.9}
A union of two connected subsets with nonempty intersection is connected; iterate.

\subsection*{Hint to Exercise VI.05.10}
Pull a hypothetical closed decomposition of the image back to the source.

\subsection*{Hint to Exercise VI.05.11}
Use the open-set criterion for irreducibility.

\subsection*{Hint to Exercise VI.05.12}
Use domains for irreducibility and idempotents/product decompositions for connectedness.

\subsection*{Hint to Exercise VI.05.13}
For the reverse inclusion in \((x)\cap(y,z)=(xy,xz)\), write an element of \((y,z)\) as
\(ay+bz\) and use divisibility by \(x\).

\subsection*{Hint to Exercise VI.05.14}
The three ideals \((x)\), \((x-1)\), and \((x-a)\) are pairwise comaximal.  Apply CRT.

\subsection*{Hint to Exercise VI.05.15}
If the image were a separation, its two inverse images would separate the source.

\subsection*{Hint to Exercise VI.05.16}
Apply \(V\) to the intersection and remember that ideal inclusion reverses zero-set
inclusion.

\subsection*{Hint to Exercise VI.05.17}
The components are \(V(x,y)\), \(V(x,z)\), and \(V(z-1,y-x^2)\).  Intersect their defining
ideals pairwise.

\subsection*{Hint to Exercise VI.05.18}
Use induction and the fact that every nonempty open subset of an irreducible space is itself
irreducible and dense.  For principal opens, use \(D(f_1)\cap\cdots\cap D(f_r)=D(f_1\cdots f_r)\).
